%%%%%%%%%%%%%%%%%%%%%%%%%%%%%%%%%%%%%%%%%%%%%%%%%%%%%%%%%%%%%%%%%%%%%%%%%%%%%%%%
% BLOQUE 12 — DOCUMENTO ADJUNTO: CUESTIONARIO DE 30 PREGUNTAS
% Requisitos en el punto 3 de las indicaciones del curso.
%
% Reinicia la numeración de página con prefijo "B" (B1, B2, ...) para
% distinguirla del Anexo de declaración de IA (prefijo "A", bloque 11).
%%%%%%%%%%%%%%%%%%%%%%%%%%%%%%%%%%%%%%%%%%%%%%%%%%%%%%%%%%%%%%%%%%%%%%%%%%%%%%%%

\pagenumbering{arabic}
\setcounter{page}{1}
\renewcommand{\thepage}{\hspace*{1em}B\arabic{page}}

\phantomsection
\addcontentsline{toc}{section}{Documento adjunto: Cuestionario de evaluación de conocimientos}
\section*{Documento adjunto: Cuestionario de evaluación de conocimientos}

% TODO (equipo) — requisitos obligatorios del cuestionario (punto 3 de
% las indicaciones):
%   - 30 preguntas en total sobre aspectos clave de las tecnologías
%     investigadas.
%   - Combinación de formatos: selección múltiple, verdadero/falso,
%     completar y respuesta corta.
%   - Cada pregunta incluye su respuesta correcta y una breve
%     justificación.
%   - Cubrir tanto teoría como aplicaciones prácticas.
%   - Dificultad graduada: 40% básicas, 40% intermedias, 20% avanzadas.
%   - Incluir preguntas que evalúen comprensión, aplicación y análisis
%     crítico (no solo memorización).
%   - Autoría: la redacción de las preguntas y justificaciones es
%     obligatoriamente humana (punto 6.1); no se admite generación por
%     IA en esta sección.
%
% Sugerencia de organización: agrupar las preguntas por sección de
% origen (ver índice temático abajo) y marcar el nivel de dificultad de
% cada una junto a su enunciado.

\subsection*{Índice temático}

% TODO (equipo): completar cuántas preguntas corresponden a cada una de
% las 10 secciones del informe y su rango (ej. "1-3"), de modo que
% sumen 30.

{\scriptsize
\begin{longtable}{|p{7.5cm}|p{2.5cm}|p{2.5cm}|}
    \hline
    \textbf{Sección del informe} & \textbf{N.\textordmasculine\ de preguntas} & \textbf{Rango} \\
    \hline
    \endhead
    1. Portada y formalidades & & \\ \hline
    2. Resumen ejecutivo & & \\ \hline
    3. Introducción y contexto & & \\ \hline
    4. Marco teórico y conceptual & & \\ \hline
    5. Desarrollo y análisis técnico & & \\ \hline
    6. Entregables específicos & & \\ \hline
    7. Vínculo con la propuesta técnico-económica & & \\ \hline
    8. Tendencias y evolución futura & & \\ \hline
    9. Conclusiones y recomendaciones & & \\ \hline
    10. Bibliografía & & \\ \hline
    \textbf{Total} & & 30 \\ \hline
\end{longtable}
}

\subsection*{Preguntas}

% TODO (equipo): redactar aquí las 30 preguntas con su respuesta y
% justificación. No hay una macro predefinida para esto en la clase;
% se sugiere un formato simple, por ejemplo:
%
%   \textbf{P1 (básica — Introducción y contexto).} Enunciado de la
%   pregunta. \textit{Respuesta:} X. \textit{Justificación:} breve
%   explicación.

\clearpage
